\chapter{Common Error Messages}
\index{error messages}\index{troubleshooting}
\label{app:errors}

This appendix lists the most frequent error messages you will encounter while learning Kaappi, along with their causes and fixes. Every message shown here is the actual output of the interpreter. Errors fall into three categories---read-time, compile-time, and runtime---plus library and FFI errors, and every diagnostic carries a stable \texttt{KP} code you can look up.

\section{How Kaappi Reports Errors}
\index{error messages!format}\index{diagnostic codes}

When an error occurs while running a file, Kaappi prints the source location, a diagnostic code, the message, and the offending line:

\begin{lstlisting}[style=repl]
program.scm:1:1: error[KP3001]: undefined variable 'dispaly'. Did you mean 'display'?
    (dispaly "hi")
\end{lstlisting}

Four things to know about this format:

\begin{itemize}
    \item \textbf{Location prefix.} \texttt{program.scm:1:1:} names the file, line, and column of the offending expression---down to the exact sub-form---and the source line is echoed beneath the message. When the failure happens inside a procedure, a \texttt{called from} line points at the call site. In the REPL the location prefix is omitted.
    \item \textbf{Diagnostic codes.} The \texttt{[KP3001]} tag identifies the diagnostic independently of its wording: \texttt{KP1xxx} are read-time, \texttt{KP2xxx} compile-time and library, \texttt{KP3xxx} runtime, and \texttt{KP4xxx} the lint warnings of \texttt{kaappi check}. Run \texttt{kaappi explain KP3001}\index{kaappi explain} for any code's meaning, a minimal example, and the fix; programs can dispatch on codes with \texttt{error-object-code} (Section~\ref{sec:error-codes}).
    \item \textbf{Did-you-mean suggestions.} For misspelled variables, Kaappi suggests the closest known binding. The suggestion depends on what is in scope, so this appendix quotes only the stable part of each message.
    \item \textbf{Machine-readable output.} \texttt{--diagnostics=json} emits every diagnostic as LSP-shaped JSON on stderr for editors and scripts (Appendix~\ref{app:cli}).
\end{itemize}

A raised exception that reaches the top level is reported with the error object's message and irritants, e.g.\ \texttt{error[KP3004]: division by zero}. Handlers can also read both fields programmatically with \code{guard}:

\begin{lstlisting}[style=repl]
> (guard (e (#t (display (error-object-message e))))
    (/ 1 0))
division by zero
\end{lstlisting}

\section{Read-Time Errors}
\index{read error}

The reader reports these when it cannot parse your source into Scheme data, with a \texttt{file:line:col:} prefix.

\subsection{\texttt{read error[KP1001]: unexpected end of input}}

\textbf{Cause:} Input ended while a list, string, or block comment was still open. Usually a missing closing parenthesis.

\begin{lstlisting}[style=repl]
(define (square x)
  (* x x)
; program.scm:3:1: read error[KP1001]:
;   unexpected end of input
\end{lstlisting}

\textbf{Fix:} Add the missing \texttt{)}. Use an editor with bracket matching; the REPL's multi-line mode tracks open parens for you.

\subsection{\texttt{read error[KP1003]: unexpected ')'}}

\textbf{Cause:} More closing \texttt{)} than opening \texttt{(}.

\begin{lstlisting}[style=repl]
(display (+ 1 2)))
; program.scm:1:19: read error[KP1003]:
;   unexpected ')'
\end{lstlisting}

\textbf{Fix:} Remove the extra parenthesis.

\subsection{\texttt{read error[KP1006]: unterminated string literal}}

\textbf{Cause:} A string literal is missing its closing double quote.

\textbf{Fix:} Add the closing \texttt{"}, or check for unescaped quotes inside the string. Use \texttt{\textbackslash"} to include a literal quote.

\subsection{\texttt{read error[KP1007]: invalid escape sequence}}

\textbf{Cause:} A string contains an escape sequence Kaappi does not recognize, such as \texttt{\textbackslash q}. Valid escapes are \texttt{\textbackslash n}, \texttt{\textbackslash r}, \texttt{\textbackslash t}, \texttt{\textbackslash a}, \texttt{\textbackslash b}, \texttt{\textbackslash\textbackslash}, \texttt{\textbackslash"}, \texttt{\textbackslash|}, and \texttt{\textbackslash x<hex>;}.

\textbf{Fix:} Use a supported escape. For hex escapes, write hex digits followed by a semicolon: \texttt{\textbackslash x41;}.

\section{Compile-Time Errors}

\subsection{\texttt{compile error[KP2001]: invalid syntax}}

\textbf{Cause:} A special form is malformed---\texttt{let} with no bindings or body, \texttt{lambda} with no parameter list, \texttt{if} with no branches, \texttt{define} with no value.

\begin{lstlisting}[style=repl]
> (lambda)
; compile error[KP2001]: invalid syntax
\end{lstlisting}

\textbf{Fix:} Supply all required sub-expressions. \texttt{lambda} needs at least a parameter list and a body: \code{(lambda (x) x)}.

\section{Runtime Errors}

\subsection{\texttt{error[KP3001]: undefined variable 'foo'}}

\textbf{Cause:} You referenced a name that has no binding. When a similarly spelled binding exists, Kaappi appends a suggestion: \texttt{Did you mean 'log'?}

\textbf{Common reasons:}
\begin{itemize}
    \item Typo in the variable name (e.g., \texttt{lenght} instead of \texttt{length})---follow the did-you-mean hint.
    \item Used a name before its \texttt{define} (order matters in scripts).
    \item The name comes from a \emph{portable} SRFI or an ecosystem library you have not imported. Kaappi's eight built-in SRFIs (including SRFI-1's \texttt{filter} and \texttt{fold}) are always available, but the 64 portable SRFIs load from \texttt{.sld} files only when imported---\texttt{list-sort} needs \texttt{(import (srfi 132))}. Appendix~\ref{app:srfi} marks which SRFIs are built in.
\end{itemize}

\textbf{Fix:} Check the spelling, reorder definitions, or add the missing \texttt{(import ...)}.

The \texttt{set!} variant of this error reads \texttt{error[KP3001]: set!: unbound variable 'y'}---\texttt{set!} requires an existing binding, so use \code{define} for the initial one.

\subsection{\texttt{error[KP3002]: type error in 'car': expected pair, got ()}}

\textbf{Cause:} A primitive received an argument of the wrong type. The message names the procedure, the expected type, and the offending argument---here, \texttt{car} applied to the empty list. Some primitives, such as the arithmetic operators, show the argument's type rather than its value:

\begin{lstlisting}[style=repl]
> (car '())
; error[KP3002]: type error in 'car':
;   expected pair, got ()
> (+ 1 "two")
; error[KP3002]: type error in 'arithmetic':
;   expected number, got #<string>
\end{lstlisting}

\textbf{Fix:} Check for the empty list with \code{null?} before calling \code{car}/\code{cdr}. For arithmetic on strings, convert first with \code{string->number}.

\subsection{\texttt{error[KP3003]: 'f': expected 2 arguments, got 1}}

\textbf{Cause:} A procedure was called with the wrong number of arguments. Variadic procedures report \texttt{expected at least N arguments}.

\textbf{Fix:} Check the procedure's signature. Use \texttt{,describe f} in the REPL to see its arity. \texttt{kaappi check} catches wrong-arity calls to built-ins before the program runs.

\subsection{\texttt{error[KP3005]: not a procedure}}

\textbf{Cause:} The first element of a call expression evaluated to something that is not callable: \texttt{(42 1)}, or \texttt{(x y)} where \texttt{x} holds a number.

\textbf{Fix:} Check for extra parentheses---\code{((+ 1 2))} computes \texttt{3} and then tries to call it.

\subsection{\texttt{error[KP3004]: division by zero} (uncaught exceptions)}

\textbf{Cause:} An exception raised by \code{error}, \code{raise}, or a primitive (such as division by zero) reached the top level uncaught. The report shows the error object's message followed by its irritants, if any. Exceptions the runtime raises carry their own specific code, like \texttt{KP3004} here; error objects your code creates with \code{error} are reported under the generic \texttt{KP3000}:

\begin{lstlisting}[style=repl]
> (/ 1 0)
; error[KP3004]: division by zero
> (error "bad input" 42)
; error[KP3000]: bad input 42
\end{lstlisting}

\textbf{Fix:} Wrap the code in \code{guard} to catch the exception and handle it (see Chapter~\ref{ch:errors}). To avoid dividing by zero in the first place, test the divisor: \code{(if (zero? b) ... (/ a b))}.

\subsection{\texttt{error[KP3006]: vector-ref: index 5 out of range for length 1}}

\textbf{Cause:} \texttt{vector-ref}, \texttt{string-ref},
or \texttt{bytevector-u8-ref} called with an index beyond
the length. The message names the procedure, the
out-of-range index, and the sequence length.

\textbf{Fix:} Check that the index is between 0 and
the sequence length minus one. Use the matching
length procedure to verify bounds.

\subsection{\texttt{error[KP3008]: stack overflow}}

\textbf{Cause:} Non-tail-recursive function with very deep recursion. The call stack grows automatically up to 32768 frames; only recursion deeper than that overflows.

\textbf{Fix:} Rewrite as a tail-recursive function using an accumulator or named \texttt{let}. See Chapter~\ref{ch:functions} for the tail-call pattern.

\section{Library and Import Errors}

\subsection{\texttt{error[KP2001]: library not found: (mylib.utils)}}

\textbf{Cause:} Kaappi cannot find the file \texttt{mylib/utils.sld} on any search path. (The message prints the library name with dots.)

\textbf{Fix:} Ensure the \texttt{.sld} file exists at the correct path relative to your working directory, \texttt{./lib/}, or a \texttt{--lib-path} directory. Ecosystem libraries are installed with \texttt{thottam install} and auto-discovered from \texttt{\textasciitilde/.kaappi/lib/}. See Appendix~\ref{app:cli} for the full search order.

\subsection{\texttt{error[KP2001]: circular import: (mylib.a)}}

\textbf{Cause:} Library A imports B, and B imports A.

\textbf{Fix:} Extract shared definitions into a third library C that both A and B import.

\section{FFI Errors}

\subsection{\texttt{error[KP3002]: ffi-open: dlopen(...)}}

\textbf{Cause:} \texttt{ffi-open} cannot find or load the shared library file (\texttt{.dylib} on macOS, \texttt{.so} on Linux). The message passes along the system loader's own report---\texttt{dlopen} on macOS lists every path it tried.

\textbf{Fix:} Build the native library first (\texttt{make} in its directory), then use the full path, pass \texttt{--lib-path}, or set \texttt{DYLD\_LIBRARY\_PATH} (macOS) / \texttt{LD\_LIBRARY\_PATH} (Linux).

\subsection{\texttt{error[KP3002]: ffi-fn: symbol not found}}

\textbf{Cause:} The shared library loaded, but does not contain the requested function. The message includes the loader's own report (\texttt{dlsym} on macOS).

\textbf{Fix:} Check the function name for typos. List the library's exported symbols with \texttt{nm -gU lib.dylib} (macOS) or \texttt{nm -D lib.so} (Linux).

\subsection{\texttt{error[KP2001]: sandbox: cannot load library from file}}

\textbf{Cause:} Running with the \texttt{--sandbox} flag, which blocks loading libraries from files---including \texttt{(kaappi ffi)} itself, so sandboxed programs cannot reach the FFI at all (file I/O, \texttt{eval}, and more are also disabled; see Appendix~\ref{app:cli}).

\textbf{Fix:} Remove \texttt{--sandbox} if you need FFI access.

\section{Pitfalls That Are Not Errors}
\index{pitfalls}

Behavior that silently surprises programmers coming from Python or JavaScript:

\begin{itemize}
    \item \textbf{Only \code{#f} is false.} \texttt{0}, \texttt{""}, and \texttt{'()} are all true in conditionals.
    \item \textbf{\code{eq?} does not compare numbers reliably.} Use \code{=} for numbers, \code{equal?} for structures. See Section~\ref{sec:equality}.
    \item \textbf{\code{set!} changes the binding, not the caller's variable.} Assigning to a parameter inside a procedure does not affect the argument at the call site.
    \item \textbf{String literals are immutable.} \code{string-set!} needs a mutable string; make one with \code{string-copy}.
    \item \textbf{\code{map} stops at the shortest list.} With lists of different lengths, extra elements are silently ignored.
\end{itemize}

\section{REPL Tips for Debugging}

\begin{center}
\small
\begin{tabularx}{\textwidth}{@{}lY@{}}
\textbf{Command} & \textbf{Purpose} \\
\hline
\texttt{,type expr} & Show the type of a value \\
\texttt{,describe name} & Show procedure arity and type \\
\texttt{,expand expr} & Show macro expansion \\
\texttt{,dis expr} & Disassemble a procedure \\
\texttt{,break name} & Set a breakpoint on a function \\
\texttt{,step expr} & Evaluate with single-stepping \\
\end{tabularx}
\end{center}

Breakpoints drop you into the interactive debugger; its command set is listed in Appendix~\ref{app:cli}. Outside the REPL, \texttt{kaappi check file.scm} reports compile-time problems and lint findings without running the program, and \texttt{kaappi explain CODE} documents any diagnostic this appendix does not cover.
